% Source-aligned English for printed p.219.
\clearpage
\SourcePage{219}

\LayerTarget{tohoku:example:holomorphic-equivariant-sheaves}
\textsc{An example.}\quad Suppose that $G$ is a group of
automorphisms of a holomorphic variety, or more generally of a ``holomorphic
space'' $X$, with ${\mathbf O}$ the sheaf of germs of holomorphic functions on $X$.
\LayerTarget{tohoku:classification:coherent-extension-classes}
If $A$ and $B$ are two coherent analytic sheaves on which $G$ acts, every
analytic sheaf that is an extension of $A$ by $B$ is coherent. It follows that
the classes of coherent $G$-${\mathbf O}$-sheaves that are extensions of $A$ by $B$
correspond to the elements of the complex vector space
$\operatorname{Ext}_{{\mathrm O},G}^{1}(X;A,B)$.
\LayerTarget{tohoku:equivalence:holomorphic-bundles-locally-free-sheaves}
On the other hand, if $n$ is an integer $>0$, there is a natural
correspondence between holomorphic vector bundles $E$ with fiber
$\mathbf C^n$ and coherent analytic sheaves $M$ on $X$ that are locally
isomorphic to ${\mathbf O}^n$: to $E$ corresponds the sheaf ${\mathbf O}(E)$ of germs of
holomorphic sections of $E$, and to $M$ corresponds the holomorphic bundle
whose fiber at $x\in X$ is $M(x)\otimes_{{\mathbf O}(x)}\mathbf C$ ($\mathbf C$ being
regarded as a module over the augmented algebra ${\mathbf O}(x)$).
\LayerTarget{tohoku:equivalence:equivariant-holomorphic-bundles}
It is then immediate that, under this correspondence, holomorphic vector
bundles with fiber $\mathbf C^n$ admitting $G$ as a group of operators
correspond to $G$-${\mathbf O}$-modules that, as ${\mathbf O}$-modules, are locally isomorphic to
${\mathbf O}^n$.
\LayerTarget{tohoku:classification:equivariant-bundle-extensions}
It follows that the classes of extensions of an equivariant holomorphic
bundle $E$ by another such bundle $F$ correspond naturally to the elements of
$\operatorname{Ext}_{{\mathrm O},G}^{1}(X;{\mathbf O}(E),{\mathbf O}(F))$, which, by Corollary 1 to Theorem
5.6.3, is isomorphic to the vector space
\LayerTarget{tohoku:equation:holomorphic-equivariant-ext-h1}
$H^1\bigl(X;G,{\mathbf O}(E'\otimes F)\bigr)$
($E'$ denoting the dual bundle of $E$).
\LayerTarget{tohoku:descent:discontinuous-action-characteristic-zero}
One can go further if $G$ is a \textit{discontinuous} group, for since we are
in characteristic 0, Corollary 1 to Theorem 5.3.1 gives
$H^1(X;G,L)=H^1(Y,L^G)$ for every $G$-sheaf of vector spaces. Moreover, if
$L$ is a coherent analytic sheaf, then so is $L^G$ (see [5, Chap. XII]); in
particular ${\mathbf O}(E'\otimes F)^G$ is a coherent analytic sheaf on $Y$.
\LayerTarget{tohoku:finiteness:compact-quotient-extension-classes}
It follows, for example, from [7] that if $Y$ is compact, then the vector
space of classes of equivariant holomorphic vector bundles that are
extensions of $E$ by $F$ is finite-dimensional. (In fact, spectral sequence I
shows that if $Y$ is compact and $A$ and $B$ are two coherent analytic
sheaves with operators, then the $\operatorname{Ext}_{{\mathbf O},G}^{i}(X;A,B)$ are
finite-dimensional.)
\LayerTarget{tohoku:analogue:abstract-algebraic-geometry}
There are analogous results in abstract Algebraic Geometry, except that in
that case one can write $H^*(X;G,L)=H^*(Y,L^G)$ only if, for example, the order
of $G$ is prime to the characteristic, or if $G$ acts ``without fixed
points.''

\LayerTarget{tohoku:section:5.7}
\textbf{5.7.\quad Introduction of the families $\Phi$.}
\LayerTarget{tohoku:definition:support-antifilter}
Let $\Phi$ be an antifilter of closed subsets of the space $X$, and consider
the functor

\LayerTarget{tohoku:equation:5.7.1}
\begin{equation}
\Gamma_\Phi^G(A)=\Gamma_\Phi(A)^G
\tag{5.7.1}
\end{equation}
on the category ${\mathbf C}^{X(G)}$ of abelian $G$-sheaves on $X$. Put

\LayerTarget{tohoku:equation:5.7.2}
\begin{equation}
H_\Phi^p(X;G,A)=R^p\Gamma_\Phi^G(A).
\tag{5.7.2}
\end{equation}
\LayerTarget{tohoku:construction:invariant-support-antifilter}
If $\Phi$ is the antifilter consisting of all closed subsets of $X$, one
recovers the functors $H^n(X;G,A)$ of 5.2. Let $\Phi'$ be the set of closed
subsets of $X$ contained in a closed subset $F\in\Phi$ invariant under $G$.
It is then clear that $\Phi'$ is an antifilter of closed subsets of $X$ and
that $\Gamma_\Phi^G=\Gamma_{\Phi'}^G$. We shall therefore assume henceforth
that $\Phi=\Phi'$, or equivalently that $\Phi$ \textit{is the set of closed subsets
$F$ of $X$ such that $\overline{f(F)}\in\Psi$, where $\Psi$ is an antifilter
of closed subsets of $Y$.}
\LayerTarget{tohoku:transition:section-5.7-functorial-formulas}
One then has the functorial formulas:

\TohokuContinuation{220}
